% Guía de uso: conciliación bancaria
% Compilar con XeLaTeX:   latexmk -xelatex GUIA_DE_USO.tex
\documentclass[11pt,a4paper]{article}

\usepackage[a4paper,margin=2.4cm,top=2.6cm,bottom=2.4cm]{geometry}
\usepackage{fontspec}
\usepackage{polyglossia}
\setdefaultlanguage{spanish}
\setmainfont{Liberation Serif}
\setsansfont{Liberation Sans}
\setmonofont{Liberation Mono}[Scale=0.9]

\usepackage{xcolor}
\usepackage{microtype}
\usepackage{ragged2e}
\usepackage{booktabs}
\usepackage{tabularx}
\usepackage{array}
\usepackage{colortbl}
\usepackage{enumitem}
\usepackage{titlesec}
\usepackage{fancyhdr}
\usepackage[most]{tcolorbox}
\usepackage[hidelinks]{hyperref}

% --- Colores (la misma línea que la aplicación: Bootstrap) ---
\definecolor{azul}{HTML}{0D6EFD}
\definecolor{oscuro}{HTML}{212529}
\definecolor{gris}{HTML}{6C757D}
\definecolor{grisclaro}{HTML}{F1F3F5}
\definecolor{verde}{HTML}{198754}
\definecolor{ambar}{HTML}{B58100}
\definecolor{rojo}{HTML}{DC3545}

\hypersetup{pdftitle={Guía de uso: conciliación bancaria},pdfauthor={JT Automatización},colorlinks=false}

% Sin patrones de silabeo para español en este equipo: se alinea a la izquierda para evitar espacios raros.
\RaggedRight
\setlength{\parindent}{0pt}
\setlength{\parskip}{0.55em}
\setlist{itemsep=0.25em,topsep=0.3em,parsep=0pt}
\renewcommand{\arraystretch}{1.25}
\setlength{\emergencystretch}{2em}

% --- Títulos ---
\titleformat{\section}{\sffamily\Large\bfseries\color{azul}}{\thesection.}{0.6em}{}[\vspace{-0.4em}\color{azul}\titlerule]
\titleformat{\subsection}{\sffamily\large\bfseries\color{oscuro}}{\thesubsection}{0.6em}{}
\titlespacing*{\section}{0pt}{1.6em}{0.8em}
\titlespacing*{\subsection}{0pt}{1.2em}{0.4em}

% --- Encabezado y pie ---
\setlength{\headheight}{22pt}
\addtolength{\topmargin}{-6pt}
\pagestyle{fancy}
\fancyhf{}
\renewcommand{\headrulewidth}{0.3pt}
\renewcommand{\headrule}{\hbox to\headwidth{\color{gris}\leaders\hrule height \headrulewidth\hfill}}
\fancyhead[L]{\sffamily\small\color{gris}Conciliación bancaria}
\fancyhead[R]{\sffamily\small\color{gris}Guía de uso}
\fancyfoot[C]{\sffamily\small\color{gris}\thepage}

% --- Macros para escribir la interfaz ---
\newcommand{\ui}[1]{\textsf{\textbf{#1}}}                       % botón o texto de la pantalla
\newcommand{\ruta}[1]{\textsf{\textit{#1}}}                      % ruta de menú
\newcommand{\flecha}{\,$\rightarrow$\,}
\newcommand{\cod}[1]{\texttt{#1}}

\newtcolorbox{nota}[1][Nota]{enhanced,breakable,colback=grisclaro,colframe=azul,boxrule=0pt,
  leftrule=3pt,arc=2pt,left=8pt,right=8pt,top=5pt,bottom=5pt,
  before skip=0.8em,after skip=0.8em,
  title={\sffamily\bfseries\small #1},coltitle=azul,colbacktitle=grisclaro,
  attach title to upper={\par\vspace{0.1em}},fonttitle=\sffamily}
\newtcolorbox{aviso}[1][Ojo]{enhanced,breakable,colback=yellow!8,colframe=ambar,boxrule=0pt,
  leftrule=3pt,arc=2pt,left=8pt,right=8pt,top=5pt,bottom=5pt,
  before skip=0.8em,after skip=0.8em,
  title={\sffamily\bfseries\small #1},coltitle=ambar,colbacktitle=yellow!8,
  attach title to upper={\par\vspace{0.1em}},fonttitle=\sffamily}

\newcolumntype{Y}{>{\RaggedRight\arraybackslash}X}
\newcommand{\encabezado}[1]{\sffamily\bfseries #1}

\begin{document}

% ---------------- Portada ----------------
\begin{titlepage}
\thispagestyle{empty}
\vspace*{3.2cm}
{\sffamily\fontsize{34}{40}\selectfont\bfseries\color{oscuro} Conciliación\\[0.2em] bancaria}\\[0.9em]
{\sffamily\Large\color{azul} Guía de uso}\\[0.4em]
{\color{azul}\rule{5cm}{2pt}}\\[1.4em]
{\sffamily\large\color{gris}Cómo cruzar los movimientos del banco con las facturas del SII:\\
qué se cobró, qué se pagó y qué no tiene respaldo.}

\vfill
\begin{tcolorbox}[colback=grisclaro,colframe=grisclaro,arc=3pt,left=10pt,right=10pt,top=8pt,bottom=8pt]
\sffamily\small
\textbf{Contenido}\par\vspace{0.3em}
\begin{tabular}{@{}r@{\hspace{0.8em}}l@{}}
1 & Conceptos básicos\\
2 & Primera vez: configuración inicial\\
3 & Importar archivos\\
4 & Uso diario\\
5 & La bandeja de sugerencias\\
6 & Conciliar a mano\\
7 & Usuarios y permisos\\
8 & Preguntas frecuentes\\
9 & Limitaciones actuales
\end{tabular}
\end{tcolorbox}
\vspace{0.6cm}
{\sffamily\small\color{gris}\today}
\end{titlepage}

\setcounter{tocdepth}{1}
\tableofcontents
\newpage

% ---------------- 1 ----------------
\section{Conceptos básicos}

\textbf{Para qué sirve:} cruzar los \emph{movimientos de la cuenta del banco} con las \emph{facturas del SII} (ventas por cobrar y compras por pagar), y dejar claro qué falta por cobrar, qué falta por pagar y qué movimientos no tienen respaldo.

\begin{tabularx}{\textwidth}{@{}>{\sffamily\bfseries}p{3cm}Y@{}}
\toprule
\rowcolor{grisclaro}\encabezado{Término} & \encabezado{Qué es}\\
\midrule
Empresa & La razón social que opera la cuenta bancaria: la dueña de las cartolas y del RCV. \textbf{No es un cliente.} Si hay más de una razón social, se crea una empresa por cada una y se cambia entre ellas desde el menú superior.\\
Movimiento & Una fila de la cartola: un abono (entra plata) o un cargo (sale plata).\\
Documento & Una factura o nota del RCV del SII. Las \textbf{ventas} son lo que le cobras a clientes y las \textbf{compras} lo que le pagas a proveedores.\\
Contraparte & Un cliente, proveedor o persona. Se crean solas al importar el RCV.\\
Conciliar & Decir «este movimiento corresponde a esta factura» (o a una categoría si no hay factura).\\
Categoría & Destino de un movimiento sin factura: combustible, retiro del dueño, traspaso entre cuentas propias, sueldos, etc.\\
Motor & Revisa los movimientos pendientes y propone con qué factura o categoría van.\\
\bottomrule
\end{tabularx}

% ---------------- 2 ----------------
\section{Primera vez: configuración inicial}

Conviene seguir este orden. \textbf{Importa primero el RCV y configura las contrapartes, y después las cartolas.} Así el motor ya sabe qué hacer con COPEC y con las transferencias al dueño cuando llegue a los movimientos.

\begin{enumerate}
\item \textbf{Entra} con el usuario administrador.
\item \textbf{Crea la empresa.} Si todavía no hay ninguna, la aplicación lo pide al entrar: escribe el RUT y la razón social. Se crea con las categorías iniciales (traspasos, retiro del dueño, remuneraciones, comisiones, impuestos, combustible, gastos e ingresos sin documento).
\item \textbf{Importa el RCV} de compras y de ventas (sección~3).
\item \textbf{Configura las contrapartes} en \ruta{Configuración\flecha Contrapartes}:
  \begin{itemize}
  \item \textbf{Proveedores que se concilian por categoría}, como COPEC: búscalo y en «Se concilia como» elige \ui{Categoría: Combustible y peajes}. Sus cargos se conciliarán solos y sus facturas dejarán de contarse como deuda por pagar.
  \item \textbf{El dueño:} \ui{Nueva persona o contraparte}\flecha nombre completo tal como está en el banco, tipo \ui{Persona}, categoría \ui{Retiro del dueño}.
  \item \textbf{Personas a las que se transfiere seguido} (sueldos, honorarios): lo mismo, con su categoría.
  \item Escribe el nombre completo. Si la cartola dice \cod{TEF A VALENZUELA CISTERNAS LUIS HUM}, registra \cod{VALENZUELA CISTERNAS LUIS HUMBERTO}. El banco corta el texto a unos 35 caracteres y el motor reconoce el inicio.
  \end{itemize}
\item \textbf{(Opcional) Reglas} en \ruta{Configuración\flecha Reglas} para movimientos repetidos sin factura. Ejemplo: «la glosa contiene MOVISTAR» + «Cargos»\flecha \ui{Gastos sin documento}.
\item \textbf{Importa las cartolas} (sección~3). Al terminar, el motor corre solo y deja lista la bandeja.
\end{enumerate}

Los usuarios se agregan en \ruta{Configuración\flecha Usuarios} (sección~7).

% ---------------- 3 ----------------
\section{Importar archivos}

Entra a \ui{Importar} y arrastra los archivos (o usa «Elegir archivos»). Puedes subir varios a la vez.

\begin{tabularx}{\textwidth}{@{}>{\sffamily\bfseries}p{3.3cm}Y p{2.6cm}@{}}
\toprule
\rowcolor{grisclaro}\encabezado{Archivo} & \encabezado{Cómo obtenerlo} & \encabezado{Formato}\\
\midrule
Cartola & Portal empresas de BancoEstado, «Cartola Histórica de Chequera Electrónica» & Excel (\cod{.xlsx})\\
RCV de ventas y de compras & Portal del SII, Registro de Compras y Ventas, detalle del período & CSV\\
\bottomrule
\end{tabularx}

\begin{aviso}[No les cambies el nombre]
Los CSV del SII (\cod{RCV\_COMPRA\_REGISTRO\_<rut>\_<aaaamm>.csv}) traen el RUT y el período en el nombre, y la aplicación los usa para validar.
\end{aviso}

\subsection*{Qué verifica la importación}
\begin{itemize}
\item \textbf{Archivo ya subido:} se rechaza.
\item \textbf{RUT distinto al de la empresa:} se rechaza.
\item \textbf{Cartola alterada:} si el saldo no cuadra al peso en alguna fila, se rechaza.
\item \textbf{Movimientos o documentos repetidos} (por ejemplo, dos cartolas que se traslapan): se descartan solos, y el resultado dice cuántos eran nuevos.
\item \textbf{Falta una cartola:} si te saltas la N.º 32 entre la 31 y la 33, lo avisa con una advertencia amarilla.
\item \textbf{RCV de compras «Reclamado» o «No incluir»:} se rechaza. El «Pendiente» se acepta con una advertencia.
\end{itemize}

La cuenta bancaria se crea sola con la primera cartola. Las notas de crédito de \textbf{ventas} se descuentan solas de su factura.

\begin{nota}[Cada cuánto]
Sube la cartola nueva cuando quieras revisar (semanal o mensual) y el RCV cada mes. Para conciliar pagos de facturas del mes anterior, el RCV de ese mes también tiene que estar cargado: un pago de septiembre cuya factura es de agosto necesita el RCV de agosto.
\end{nota}

% ---------------- 4 ----------------
\section{Uso diario}

Después de importar:

\begin{enumerate}
\item \ui{Resumen}: muestra el avance, cuántos movimientos faltan y el saldo por cobrar y por pagar. Con \ui{Buscar coincidencias} vuelves a correr el motor.
\item \ui{Bandeja}: revisa lo que el motor no se atrevió a conciliar solo (sección~5).
\item \ruta{Movimientos\flecha Por conciliar}: lo que quedó sin propuesta, uno a uno (sección~6).
\item \ui{Documentos}: consulta qué facturas siguen pendientes de cobro o de pago.
\end{enumerate}

\subsection*{Qué hace el motor solo}
El motor concilia \textbf{sin preguntarte} lo que tiene una confianza de \textbf{90\,\% o más}: COPEC contra su categoría, tus traspasos, el retiro del dueño, lo que cubren tus reglas, y los pagos donde el nombre y el monto calzan sin dudas. El resto queda como sugerencia.

% ---------------- 5 ----------------
\section{La bandeja de sugerencias}

Cada tarjeta muestra, a la izquierda, el movimiento del banco y, a la derecha, lo que propone el motor, con el motivo y el porcentaje de confianza.

\begin{tabularx}{\textwidth}{@{}>{\sffamily\bfseries}p{4cm}Y@{}}
\toprule
\rowcolor{grisclaro}\encabezado{Botón} & \encabezado{Qué hace}\\
\midrule
Aceptar & Concilia con la propuesta marcada.\\
Elegir entre opciones & Si hay varias, marca la que corresponde con el círculo y luego \ui{Aceptar}.\\
Rechazar & Descarta la propuesta; no vuelve a aparecer para ese movimiento.\\
Otra & Abre la conciliación manual.\\
Aceptar las visibles & Acepta todas las que estás viendo de una vez (pide confirmación).\\
\bottomrule
\end{tabularx}

\subsection*{Cómo leer el porcentaje}
\begin{itemize}
\item \textcolor{verde}{\textbf{Verde (90\,\% o más):}} casi seguro. El motor lo concilia solo.
\item \textcolor{azul}{\textbf{Azul (60 a 90\,\%):}} probable; conviene mirarlo un segundo.
\item \textcolor{ambar}{\textbf{Amarillo o gris (menos de 60\,\%):}} dudoso. Suele ser «solo coincide el monto», por ejemplo un depósito sin nombre. Revísalo con cuidado.
\end{itemize}

Usa los filtros \ui{Probables} y \ui{Dudosas} para no aceptar a ciegas las dudosas con \ui{Aceptar las visibles}.

\begin{nota}[El sistema aprende]
Cada vez que concilias un pago con una factura, queda guardado qué texto del banco corresponde a qué cliente. Con el uso, hay más conciliaciones automáticas.
\end{nota}

% ---------------- 6 ----------------
\section{Conciliar a mano}

En \ui{Movimientos}, haz clic en cualquier fila. Se abre «Conciliar movimiento»:

\begin{itemize}
\item \textbf{Una factura:} marca la factura de la lista. Vienen ordenadas por monto exacto y parecido de nombre; con el buscador puedes encontrar otra por razón social, RUT o folio.
\item \textbf{Varias facturas con un solo pago:} marca varias. Cada una trae su monto a aplicar.
\item \textbf{Pago parcial o abono:} baja el monto de la factura. El movimiento queda \emph{Parcial} hasta completar.
\item \textbf{Sin factura} (comisión, retiro, traspaso): elige una \textbf{categoría} en «Sin documento».
\item \textbf{Factura más categoría:} si el movimiento es mayor que las facturas, el resto va a una categoría.
\item \textbf{Ignorar:} para un movimiento que no hay que conciliar (duplicado, de prueba). Se puede revertir.
\item \textbf{Deshacer:} si ya estaba conciliado, el botón \ui{Deshacer} en la misma ventana devuelve todo a como estaba. Nada se borra: queda registrado quién hizo qué.
\end{itemize}

El botón \ui{Conciliar} solo se activa si lo que aplicas no supera el monto disponible. Un abono se aplica a ventas y un cargo a compras.

\subsection*{Notas de crédito de compras}
El SII no informa de qué factura es la nota de crédito de \textbf{compras}. En \ruta{Documentos\flecha Compras}, abre la nota de crédito y elige a qué factura del mismo proveedor aplicarla.

% ---------------- 7 ----------------
\section{Usuarios y permisos}

Solo un administrador entra a \ruta{Configuración\flecha Usuarios}.

\begin{tabularx}{\textwidth}{@{}>{\sffamily\bfseries}p{3cm}Y@{}}
\toprule
\rowcolor{grisclaro}\encabezado{Rol} & \encabezado{Qué puede hacer}\\
\midrule
Lector & Consultar todo.\\
Operador & Lo anterior, más importar, conciliar, aceptar o rechazar sugerencias y manejar reglas y personas.\\
Administrador & Lo anterior, más cuentas, categorías, asignar categorías a contrapartes y usuarios de la empresa.\\
\bottomrule
\end{tabularx}

El \textbf{administrador general} crea usuarios nuevos y puede restablecer contraseñas. Cada usuario cambia la suya desde su nombre, arriba a la derecha\flecha \ui{Cambiar contraseña}. Mínimo 10 caracteres. La sesión dura 8 horas.

% ---------------- 8 ----------------
\section{Preguntas frecuentes}

\textbf{¿Por qué un movimiento no tiene sugerencia?} Casi siempre es una persona no registrada, o una factura que no está cargada (falta el RCV de ese mes). Registra a la persona en \ui{Contrapartes} o carga el RCV faltante, y luego usa \ui{Buscar coincidencias}.

\textbf{Cambié una categoría o agregué una regla, ¿se aplica a lo ya importado?} Sí, pero hay que correr \ui{Buscar coincidencias}. Lo ya conciliado no se toca.

\textbf{¿Qué pasa con las facturas de COPEC?} Aparecen como \emph{Por categoría}, no como pendientes. Siguen en la base de datos pero no cuentan como deuda por pagar.

\textbf{Se me conciliaron mal varios movimientos.} Ábrelos desde \ruta{Movimientos\flecha Conciliados} y usa \ui{Deshacer}. Después revisa la contraparte o la regla que causó el error.

\textbf{¿Se guardan mis archivos?} No. Se leen en el momento y solo se guardan los datos y una huella para detectar repetidos.

\textbf{Importé una cartola dos veces.} La segunda se rechaza. Si fue una cartola distinta que se traslapa, los movimientos repetidos se descartan solos.

% ---------------- 9 ----------------
\section{Limitaciones actuales}

\begin{itemize}
\item Solo \textbf{BancoEstado} (cartola Excel). Banco de Chile requiere un lector nuevo y cartolas de ejemplo.
\item Los depósitos en efectivo y con documentos no traen el nombre de quien paga: solo se pueden proponer por monto, así que quedan para revisión.
\item Las boletas de honorarios y la liquidación de Transbank aún no se importan.
\item La descarga del banco y del SII es manual: no hay conexión automática todavía.
\end{itemize}

\end{document}
